% TOP_PROBLEM: {"id": 2800601, "title": "10 Lectures and 42 Open Problems — Random Partial Discrete Fourier Transform", "category_id": 15, "category": {"id": 15, "name": "computer_science", "display_name": "Computer Science", "description": "Computational complexity, algorithms, and theoretical CS.", "slug": "computer-science", "order_index": 15, "created_at": "2026-07-31T15:26:25.671Z"}, "status": "open", "problem_number": "AMR-027-0601", "set_id": 13}
% FIRST_POSTED: 2026-10-02
% ENTRY_KIND: statement_only
% UnsolvedMath ID 2800601; source catalogue code AMR-027-0601
% !TeX program = lualatex
% Definition and sources only; this file is not an LLM solution attempt.
\documentclass[11pt,a4paper]{article}
\usepackage[margin=25mm]{geometry}
\usepackage{fontspec}
\setmainfont{lmroman10-regular.otf}[BoldFont=lmroman10-bold.otf,
 ItalicFont=lmroman10-italic.otf,BoldItalicFont=lmroman10-bolditalic.otf]
\usepackage{amsmath,amssymb,mathtools}
\usepackage{unicode-math}
\setmathfont{latinmodern-math.otf}
\AtBeginDocument{\let\setminus\smallsetminus}
\usepackage{xurl,hyperref}
\hypersetup{colorlinks=true,urlcolor=blue}
\setcounter{secnumdepth}{0}
\setlength{\parindent}{0pt}
\setlength{\parskip}{5pt}
\setlength{\emergencystretch}{3em}
\begin{document}
\section*{10 Lectures and 42 Open Problems — Random Partial Discrete Fourier Transform}\label{2800601}
{\small UnsolvedMath ID 2800601 · AMR-027-0601 · Computer Science · AMR Open Problem Lists}
\subsection{Definitions and mathematical statement}
For an integer $N\ge2$, let the unitary discrete Fourier matrix be
$F_{rj}=N^{-1/2}\exp(2\pi i rj/N)$, with $0\le r,j<N$.
Choose $R_1,\ldots,R_M$ independently and uniformly from
$\{0,\ldots,N-1\}$, allowing repetitions, and put
\[
 A_{tj}=\sqrt{N/M}\,F_{R_tj}
       =M^{-1/2}\exp(2\pi i R_tj/N).
\]
The rescaling ensures $\mathbb E(A^*A)=I_N$ and gives every column
Euclidean norm one.

For $1\le s\le N$, the event of interest is
\[
 \frac23\|x\|_2^2\le\|Ax\|_2^2\le\frac43\|x\|_2^2
 \quad\text{for every }x\in\mathbb C^N
       \text{ with }|\operatorname{supp}(x)|\le s.
\]
For $0<\eta<1/2$, determine the sharp order of the smallest sample
size $M=M(N,s,\eta)$ for which this event has probability at least
$1-\eta$. In particular, determine the necessary logarithmic factors
in the high-probability regimes, for example $\eta=N^{-1}$.

The probability concerns the sampled rows, while the assertion within
the event is uniform over every support and every coefficient vector.
Recovery of one fixed sparse signal need not require the same number
of measurements. Rows are sampled with replacement, exactly as in
the primary formulation; sampling a subset without replacement is a
related but different random model. Both lower and upper bounds are
needed to identify the sample scale rather than merely provide a
sufficient oversampling rule.
\subsection{Short English statement}
How many independently sampled Fourier rows are necessary and sufficient to preserve all sparse vectors with high probability?
\subsection{Sources}
\begin{itemize}
\item[S1] ULAM.AI and the UnsolvedMath Contributors, supplied problems.json, record 2800601 (AMR-027-0601), AMR Open Problem Lists. Catalogue identity and source transcription; CC BY 4.0.
\url{https://huggingface.co/datasets/ulamai/UnsolvedMath}.
\item[S2] Bandeira - 42 Open Problems in Data Science (2016), Open Problem 6.1. Original source indexed by AMR-027-0601.
\url{https://afonsobandeira.wordpress.com/2015/12/07/10l42pcompressedsensing/}.
\item[S3] A. S. Bandeira, Ten Lectures and Forty-Two Open Problems in the Mathematics of Data Science, corresponding numbered problem and surrounding definitions.
\url{https://people.math.ethz.ch/~abandeira/TenLecturesFortyTwoProblems.pdf}.
\end{itemize}
\end{document}
